% ECONOMICS_PROBLEM: {"id": "F34-105", "title": "Restructuring with diverse creditors", "jel_code": "F34", "source_id": "OP-0256"}
% FIRST_POSTED: 2026-10-09
% ATTEMPT_STATUS: unresolved
% ENTRY_KIND: research_attempt
\documentclass[11pt,a4paper]{article}
\usepackage[T1]{fontenc}
\usepackage[utf8]{inputenc}
\usepackage[margin=27mm]{geometry}
\usepackage{amsmath,amssymb,amsthm,mathtools}
\usepackage[hidelinks]{hyperref}
\setlength{\parskip}{5pt}
\setlength{\parindent}{0pt}
\newtheorem{theorem}{Theorem}[section]
\newtheorem{proposition}[theorem]{Proposition}
\newtheorem{lemma}[theorem]{Lemma}
\newtheorem{corollary}[theorem]{Corollary}
\theoremstyle{remark}
\newtheorem{remark}[theorem]{Remark}
\newcommand{\R}{\mathbb R}
\newcommand{\E}{\mathbb E}
\newcommand{\Prob}{\mathbb P}
\newcommand{\ind}{\mathbf 1}
\DeclareMathOperator{\Var}{Var}
\DeclareMathOperator{\KL}{KL}
\DeclareMathOperator{\TV}{TV}
\DeclareMathOperator{\OPT}{OPT}
\title{Truthful unanimity workouts with endogenous issuance prices}
\author{GPT 6 Astra Ultra}
\date{9 October 2026}
\begin{document}\maketitle
\textbf{Problem ID:} \texttt{F34-105}. \textbf{Status:} Unresolved. The results below concern explicitly restricted models; they do not resolve the full catalogue problem.

\section{Question and a restricted creditor environment}
The catalogue asks for restructuring mechanisms that accommodate diverse creditors, legal priorities, private information, and the effect of anticipated settlement on initial borrowing. We solve a threshold-family problem with public protected claims and heterogeneous private outside recoveries. The calculation includes competitive issuance prices and a terminal repayment rule. It does not claim optimality over all mechanisms or explain actual legal systems.

There are three dates: issuance at date 0, possible immediate settlement at date 1, and terminal liquidation at date 2. Interest is zero and all parties are risk neutral in money. Public protected creditors must receive total amount $s\ge0$ under either resolution; this obligation has legal priority. Date-1 capacity is $s+B$, so $B\ge0$ is the residual capacity for $n\ge2$ flexible creditors. Flexible creditor $i$ has a privately known terminal recovery $\theta_i\in[0,H_i]$, where the independent distributions are uniform and all $H_i>0$. If no unanimous settlement occurs, adjudicated liquidation at date 2 pays each its true $\theta_i$, independently of its report. The sovereign incurs an additional real delay cost $D>0$. Assume terminal resources can fund $s+\sum_iH_i$, and forbid borrowing after date 0. Face claims may exceed these capacities, so the sovereign is insolvent relative to contractual face debt. No new financing is used to hide a repayment shortfall.

The private report concerns the value of the creditor's outside option, not its power to rewrite the terminal adjudication. This distinction is essential for the incentive argument. Public creditors' priority payment is fixed; flexible creditors are heterogeneous through $H_i$. The model does not include strategic voting by the protected creditors.

\section{A truthful threshold family}
Fix
\[
 0\le k\le\bar k:=\min\{H_1,\ldots,H_n,B/n\}.
\]
For reports $\widehat\theta$, settle if $\sum_i\widehat\theta_i\le k$. On settlement pay
\[
 p_i(\widehat\theta)=k-\sum_{j\ne i}\widehat\theta_j
\]
to each flexible creditor. Otherwise implement the specified terminal outside option. The protected claims receive $s$ in either case. All settlement payments are nonnegative.

\begin{theorem}[Truthfulness, participation, and a sharp family budget bound]
For every $k\le\min_iH_i$, the threshold rule is dominant-strategy truthful and individually rational relative to terminal recovery for each flexible creditor. Its settlement payments satisfy the residual resource constraint for every report profile if and only if $nk\le B$. Thus the interval $[0,\bar k]$ is exactly the feasible parameter interval for this family.
\end{theorem}
\begin{proof}
Fix the reports of all other creditors and write $h_i=k-\sum_{j\ne i}\widehat\theta_j$. If $h_i<0$, settlement is impossible for any nonnegative report by $i$, so all its reports yield $\theta_i$. If $h_i\ge0$, a report at most $h_i$ yields settlement payment $h_i$, while a higher report yields outside recovery $\theta_i$. Truthful reporting selects settlement precisely when $\theta_i\le h_i$, which maximizes its payoff; at equality it is indifferent. Truthful settlement therefore never pays less than the outside option. This proves dominant strategies and participation without averaging other creditors' types.

If settlement occurs, letting $z=\sum_i\widehat\theta_i\le k$ gives total payment $nk-(n-1)z\le nk$. The maximum $nk$ is attained at the zero report profile, so budget feasibility is equivalent to $nk\le B$. The public senior claim is already removed from the residual budget.
\end{proof}

Unanimity is implemented by each creditor's threshold participation decision. Truthfulness here does not require observing its realized private outside value at date 1. The all-zero profile matters for an ex post budget rule even though it has probability zero under continuous types: nearby profiles approximate the same maximum.

\section{Exact ex ante calculations}
Let $H=\prod_iH_i$ and $P(k)$ be the settlement probability. Let $Q(k)$ be the total competitively priced proceeds at issuance when creditors anticipate this mechanism, including the protected claims. Competitive risk-neutral lenders pay the expected total recovery, with no hidden guarantee.

\begin{proposition}[Settlement probability, payments, and issuance]
For $0\le k\le\bar k$,
\begin{align*}
 P(k)&=\frac{k^n}{n!H},\\
 \E\left[\ind\{\text{settle}\}\sum_i p_i\right]
 &=\frac{2n k^{n+1}}{(n+1)!H},\\
 Q(k)&=Q_0+\frac{n k^{n+1}}{(n+1)!H},
 \qquad Q_0=s+\frac12\sum_iH_i.
\end{align*}
For $k>0$, $Q'(k)=kP'(k)$.
\end{proposition}
\begin{proof}
Since $k\le\min_iH_i$, the settlement region is the full simplex $\{\theta\ge0:\sum_i\theta_i\le k\}$ inside the rectangular type support. Its volume is $k^n/n!$, proved by recursively integrating the last coordinate over $[0,k-\sum_{i<n}\theta_i]$. Uniform density is $1/H$, giving $P(k)$.

The same integration, or symmetry and one additional coordinate integral, gives
\[
 \int_{\sum\theta_i\le k}\theta_i\,d\theta
 =\frac{k^{n+1}}{(n+1)!}.
\]
Total settlement payment is $nk-(n-1)\sum_i\theta_i$. Integrating it yields
\[
 \frac1H\left[\frac{nk^{n+1}}{n!}
 -\frac{n(n-1)k^{n+1}}{(n+1)!}\right]
 =\frac{2nk^{n+1}}{(n+1)!H}.
\]
Compared with always liquidating, each settling creditor receives an increment $p_i-\theta_i=k-\sum_j\theta_j$. Summing and integrating gives incremental expected recovery $nk^{n+1}/((n+1)!H)$. Add baseline expected recovery $Q_0$. Differentiate both formulas to obtain $Q'=kP'$.
\end{proof}

The increment in $Q$ is an information rent relative to liquidation, not a new resource available free of future repayment. Competitive issuance capitalizes it exactly.

\section{Anticipated borrowing changes the optimal threshold}
Let $V(Q)$ be the sovereign's date-zero value from using the proceeds, in the same resource units as later repayments and delay. Its ex ante objective within the threshold family is
\[
 \mathcal W(k)=V(Q(k))-Q(k)-D[1-P(k)]. \tag{1}
\]
This is the sovereign's objective, not a sum of sovereign and creditor surplus: promised repayments are charged once, and proceeds enter only through $V$. The $Q(k)$ term equals expected subsequent debt service by competitive pricing. Political or distributional social weights would define another objective.

\begin{theorem}[Exact optimum in the linear-investment case]
Suppose $V(Q)=\gamma Q$ with $0\le\gamma<1$ and $\bar k>0$. The unique maximizer of (1) is
\[
 k^*=\min\{\bar k,D/(1-\gamma)\}.
\]
For $\gamma\ge1$, the unique maximizer is $\bar k$. If $V$ is differentiable, an interior optimum $k>0$ must instead satisfy
\[
 D=[1-V'(Q(k))]k. \tag{2}
\]
When $V$ is concave and $V'\le1$ on the relevant proceeds interval, the derivative's sign changes at most once from positive to negative, so any root of (2) is a global maximum and the endpoints are checked when no root exists.
\end{theorem}
\begin{proof}
Differentiate (1) and use $Q'=kP'$ to obtain
\[
 \mathcal W'(k)=P'(k)\{D+[V'(Q(k))-1]k\}.
\]
For $k>0$, $P'(k)>0$. Under linear $V$ the bracket is $D-(1-\gamma)k$, proving the formulas and uniqueness, including the cases where the unconstrained root lies at or beyond the upper endpoint. Although $P'(0)=0$ for $n\ge2$, the objective increases for all sufficiently small positive $k$, so zero is not another maximizer when $\bar k>0$.

For concave $V$ with $V'\le1$, the function $k\mapsto k[1-V'(Q(k))]$ is nondecreasing: both nonnegative factors are nondecreasing, since $Q$ increases and $V'$ decreases. Therefore the bracket can cross zero only in the asserted direction. A zero interval gives multiple global maximizers rather than a false uniqueness claim. Equation (2) and endpoint checking follow.
\end{proof}

With issuance proceeds already sunk, the date-1 choice within this family maximizes $DP(k)-Q(k)$ and therefore selects $k=\min\{\bar k,D\}$. Unless $\gamma=0$ or a resource bound binds both choices, this differs from the ex ante commitment threshold. The difference is generated by borrowing terms, not by treating immediate debt relief as a free transfer.

\section{Legal and empirical limitations}
A public creditor demanding protected repayment greater than total immediate capacity makes the specified immediate regime infeasible before private incentives are considered. Conversely, changing a legal priority rule changes $s$ and $B$ and can alter $\bar k$ discretely. The threshold formulas isolate that legal channel from the private-information channel, whose observable implication is the relation between settlement probability and expected recovery.

Identification would require documented claim priorities, recovery data including unsettled cases, and a model separating private outside recoveries from unobserved sovereign capacity. Observed settlement payments alone do not identify all $H_i$ or $V'$. Correlated private types, strategic official creditors, multilateral financing, richer votes, and the best mechanism outside this threshold family remain unresolved.

\begin{thebibliography}{9}
\bibitem{MT} K. J. Mitchener and C. Trebesch, \emph{Sovereign Debt in the 21st Century}, CESifo Working Paper 8959. Original version March 2021; the primary PDF currently served is marked September 2022. Sections 8--10 discuss official lending, legal risks, and research directions. \url{https://www.ifo.de/DocDL/cesifo1_wp8959.pdf}. The mechanism and uniform-type formulas here are editorial restrictions, not claims attributed to that survey.
\end{thebibliography}

\end{document}
